\begin{longtable}{rllrp{7cm}}
\toprule
\cell{\#} & Group & \cell{Effect} & \cell{N} & Note \\
\midrule
\endhead
1 & North, treated & -1.69 & 12 & The sample of respondents measures the share of workers who switch from driving to cycling after construction, although the confidence interval is wide. \\
2 & North, treated & +2.53 & 13 & Our estimate confirms the cost of construction relative to the recorded benefits for households, given the limits of the survey instrument. \\
3 & South, treated & -2.74 & 6 & The main specification measures considerable heterogeneity between older and newer residential areas, although the confidence interval is wide. \\
4 & West, control & +4.56 & 7 & The regression model measures a modest reduction in average commuting time across the districts studied, while the effect on weekend trips is close to zero. \\
5 & East, treated & +1.94 & 13 & Our estimate illustrates the share of workers who switch from driving to cycling after construction, and the result does not depend on outliers. \\
6 & West, treated & +5.93 & 6 & Local planners compares a modest reduction in average commuting time across the districts studied, while the effect on weekend trips is close to zero. \\
7 & South, control & -2.53 & 5 & A typical resident illustrates the relationship between lane length and measured travel time over several years, although the confidence interval is wide. \\
8 & North, control & -6.88 & 20 & The sample of respondents highlights a stable pattern that survives alternative definitions of the treatment period, while the effect on weekend trips is close to zero. \\
9 & East, treated & -1.34 & 12 & The sample of respondents highlights the sensitivity of the estimates to the choice of comparison neighbourhoods, and the result does not depend on outliers. \\
10 & North, treated & +2.14 & 16 & The sample of respondents compares the sensitivity of the estimates to the choice of comparison neighbourhoods, given the limits of the survey instrument. \\
11 & West, treated & +5.59 & 13 & Our estimate suggests the share of workers who switch from driving to cycling after construction, once we control for population density. \\
12 & West, control & -7.35 & 15 & The main specification highlights a modest reduction in average commuting time across the districts studied, although the confidence interval is wide. \\
13 & South, treated & +7.60 & 11 & The baseline comparison highlights considerable heterogeneity between older and newer residential areas, although the confidence interval is wide. \\
14 & West, treated & -8.53 & 3 & Local planners illustrates considerable heterogeneity between older and newer residential areas, and the result does not depend on outliers. \\
15 & South, control & -3.82 & 6 & The pilot programme illustrates the relationship between lane length and measured travel time over several years, although the confidence interval is wide. \\
16 & East, control & -6.14 & 20 & The regression model estimates the share of workers who switch from driving to cycling after construction, once we control for population density. \\
17 & West, treated & +1.65 & 20 & The main specification compares the cost of construction relative to the recorded benefits for households, once we control for population density. \\
18 & North, treated & -5.72 & 6 & The pilot programme describes the cost of construction relative to the recorded benefits for households, while the effect on weekend trips is close to zero. \\
19 & South, control & -6.93 & 17 & The city council rejects the cost of construction relative to the recorded benefits for households, while the effect on weekend trips is close to zero. \\
20 & North, control & +8.87 & 13 & The baseline comparison reveals a stable pattern that survives alternative definitions of the treatment period, while the effect on weekend trips is close to zero. \\
\bottomrule
\end{longtable}
